\documentclass[a4paper,2pt]{article}
\usepackage{fontspec}
\usepackage[utf8]{inputenc}
\setmainfont{Times New Roman} % Replace with a Unicode-compatible font installed on your system
\usepackage[russian]{babel}
\usepackage{geometry}
\usepackage{graphicx}
\usepackage{array}
\usepackage{amsmath, amssymb}
\geometry{left=0cm,right=1cm,top=0cm,bottom=0cm}
\begin{document}
\thispagestyle{empty}
\fontsize{4pt}{3pt}\selectfont
% {\tiny

\begin{center}
  {\large\bf Шпаргалка кр1}
\end{center}

\vspace{0.5em}

\setlength{\extrarowheight}{1pt} % Reduce row height
\begin{tabular}{|p{0.315\textwidth}|p{0.315\textwidth}|p{0.315\textwidth}|} % Adjust column widths
\hline

\textbf{Конечно-разностный метод для решения краевых задач ОДУ.:}
Решение краевых задач заключается в их сведении к решению систем линейных и нелинейных уравнений. Одними из наиболее популярных являются конечно-разностные методы.
$\frac{d^2y}{dt^2} + p(t)\frac{dy}{dt} + q(t)y(t) = f(t)$$\alpha_1 y(a) + \beta_1 \frac{dy(a)}{dt} = \gamma_1\qquad \alpha_2 y(b) + \beta_2 \frac{dy(b)}{dt} = \gamma_2$

Введем дискретные значения независимой переменной. Аппроксимируем первую и вторую производную, и в итоге получаем систему относительно $y_k$
$a_ky_{k-1}+b_ky_k+c_ky_{k+1}=-h^2f_k \tag{1}$

Граничные условия принимают вид:
$\alpha_1y_0+\beta_1\frac{-y_2+4y_1-3y_0}{2h}=\gamma_1\qquad \alpha_2y_N+\beta_2\frac{3y_N-4y_{N-1}+y_{N-2}}{2h}=\gamma_2 \tag{2}$

(2) подставляем в (1), получаем
$\tilde{b}_1y_1+\tilde{c}_1y_2=-h^2\tilde{f}_1\qquad  \tilde{a}_{N-1}y_{N-2}+\tilde{b}_{N-1}y_{N-1}=-h^2\tilde{f}_{N-1} \tag{3}$

(1) и (3) порождают линейную систему с трехдиагональной матрицей

Метод прогонки:
Решение (1) и (3) ищем в следующем виде
$y_k=\Theta_ky_{k+1}+\Psi_k \qquad(4)$ $y_{k-1}=\Theta_{k-1}y_k+\Psi_{k-1}=\Theta_{k-1}\Theta_ky_{k+1}+(\Theta_{k-1}\Psi_k+\Psi_{k-1}) \tag{5}$
Подставим (4) и (5) в (1)
$[\Theta_k(a_k\Theta_{k-1}+b_k)+c_k]y_{k+1}+[\Psi_k(a_k\Theta_{k-1}+b_k)+a_k\Psi_{k-1}+h^2f_k]=0$

Прямая прогонка
$\Theta_k=-\frac{c_k}{b_k+a_k\Theta_{k-1}} \qquad \Psi_k=-\frac{a_k\Psi_{k-1}+h^2f_k}{b_k+a_k\Theta_{k-1}}$  $\Theta_1=-\frac{\tilde{c}_1}{\tilde{b}_1} \qquad \Psi_1=-\frac{h^2\tilde{f}_1}{\tilde{b}_1}  \tag{6}$

Обратная прогонка
$y_{N-1}=-\frac{\tilde{a}_{N-1}\Psi_{N-2}+h^2\tilde{f}_{N-1}}{\tilde{b}_{N-1}+\tilde{a}_{N-1}\Theta_{N-2}}  \tag{7}$


&

\textbf{Методы коллокации и Галеркина для решения краевых задач:}
задача:
$\frac{d}{dt}(p(t)\frac{d}{dt}y(t)) + q(t)y(t) = f(t)$, краевые условия: $y(0)=y(1)=0$
Приближенное решение найдем в виде:
$y(t) = \sum_{k=1}^nc_k\phi_k(t)$

1) Метод коллокации
Ac = f, $\sum_{k=1}^na_{mk}c_k = f(t_m)$
$a_{mk} = p(t_m)\phi''_k(t_m) + p'(t_m)\phi'_k(t_m)+q(t_m)+\phi_k(t_m)$

2) Метод Галеркина
Ac = f,  $\sum_{k=1}^n a_{mk}c_k = \int_0^1f(t)\phi_m(t)$
$a_{mk} = -\int_0^1 p(t)\frac{d\phi_k(t)}{dt}\frac{d\phi_m(t)}{dt}dt$  +  $\int_0^1q(t)\phi_k(t)\phi_m(t)dt$

Примеры базисных функций, удовлетворяющих краевым условиям:
1. $\phi_k(t)=sin(k\pi)$
2. $\phi_k(t)=t^k(1-t)$

Два замечания по поводу систем уравнений: 1. Матрицы коэффициентов в обоих методах в общем случае являются заполненными. 2. Преимущество метода Галеркина перед методом коллокации в задаче выше в том, что его матрица симметрична.


&

\textbf{Примеры уравнений в частных производных. Метод сеток и метод прямых:}
Примеры уравнений:
1. Уравнение переноса: $\frac{\partial u}{\partial t} = -a\frac{\partial u}{\partial x}+f(t,x)$
2. Уравнение теплопроводности (параболического типа): $\frac{\partial u}{\partial t} = a(\frac{\partial^2 u}{\partial x^2}+\frac{\partial^2 u}{\partial y^2}+ \frac{\partial^2 u}{\partial z^2})$
3. Волновое уравнение (гиперболического типа): $\frac{\partial^2 u}{\partial t^2} = a(\frac{\partial^2 u}{\partial x^2}+\frac{\partial^2 u}{\partial y^2}+ \frac{\partial^2 u}{\partial z^2})$
4. Уравнение потенциала (уравнения Пуассона, Лапласа, эллиптического типа): $\frac{\partial^2 u}{\partial x^2}+\frac{\partial^2 u}{\partial y^2}+ \frac{\partial^2 u}{\partial z^2} = f(x,y,z)$
Для выделения одного решения необходимо задать начальные и/или краевые условия:
- условие Дирихле, когда задается начальное распределение функции u(t,x,y,z) на границах области ее изменения
- условие Неймана, когда определяется изменение функции u(t,x,y,z) на границе.

1) Метод сеток
Задача: $a\frac{\partial^2 u}{\partial x^2}+b\frac{\partial^2 u}{\partial y^2}+ c\frac{\partial u}{\partial x} + d\frac{\partial u}{\partial y}+eu = f(x,y)$

Введем дискретизацию по обеим независимым переменным.
Граничные условия:
$u(A, y_i) = \phi_1(y_i)$, $u(B, y_i) = \phi_2(y_i)$, $u(x_k, C) = \phi_3(x_k)$, $u(x_k, D) = \phi_4(x_k)$

Запишем исходное равнение в точке $(x_k,y_i)$, приближенно аппроксимируя производные по формулам численного дифференцирования:
$\frac{\partial^2u}{\partial x^2}(x_k, y_i)\simeq\frac{u_{k+1,i} - 2u_{k,i} + u_{k-1,i}}{h^2_x}$; $\frac{\partial^2u}{\partial y^2}(x_k, y_i)\simeq\frac{u_{k,i+1} - 2u_{k,i} + u_{k,i-1}}{h^2_y}$,
$\frac{\partial u}{\partial x}(x_k, y_i)\simeq\frac{u_{k+1,i} - u_{k-1,i}}{2h_x}$; $\frac{\partial u}{\partial y}(x_k, y_i)\simeq\frac{u_{k,i+1} - u_{k,i-1}}{2h_y}$

Количество уравнений (N-1)(M-1). Оставшиеся значения $u_{k,i}$ известны и отвечают границам прямоугольника.
Если область на плоскости имеет относительно произвольный вид, то уравнения можно записать только в тех "внутренних" точках $(x_k, y_i)$, которые имеют четырех соседей. Для точек, расположенных внутри области, но не имеющих хотя бы 1 "соседа", реализуют снос граничных условий, задавая этим $u_k,i$ значения ближайших точек границы области. Повышение точности решения достигается увеличением значений N и M.

2) Метод прямых
 Заключается в дискретизации по всем независимым переменным, кроме одной.
 Задача: $\frac{\partial u(x,t)}{\partial t} = a(x,t)\frac{\partial^2 u(x,t)}{\partial x^2}$
 Введем дискретизацию. Используем формулу численного дифференцирования для второй производной, получим систему дифференциальных уравнений относительно $u_k(t)$:
 $\frac{du_k(t)}{dt} = \frac{a_k(t)}{(\Delta x)^2}(u_{k+1}(t) - 2u_k(t) + u_{k-1}(t))$

Уменьшение величины $\Delta x$ повышает точность аппроксимации, но система уравнений становится более жесткой.

\\\hline


\textbf{Минимизация функции одной переменной. Методы половинного деления:}
Предположим, что для функции f(x), определенной на \([a, b]\) , имеется единственное значение, отвечающее минимуму f на \([а ,b]\). При этом для значений \(x>x*\) функция строго возрастает и для \(x<x*\) строго убывает. Такая функция называется унимодальной. В дальнейшем речь идет об унимодальных на заданном промежутке функциях.
Возможны три варианта расположения минимума и сокращения промежутка неопределенности: 1) f(x0)>f(x1); 2) f(x0)<f(x1); 3) f(x0)=f(x1)

1) Метод половинного деления
На очередном шаге вычисляется функция в двух точках:
$x_0 = \frac{a+b}{2} - \frac{\varepsilon}{2}$ и $x_1 = \frac{a+b}{2} + \frac{\varepsilon}{2}$ с точностью до $\frac{\varepsilon}{2}$

2) Метод золотого сечения
Расстояние от 0 до $x_{лев}$ и от $x_{прав}$ до 1 обозначим $\alpha$, тогда между $x_{лев}$ и $x_{прав}$ $1-2\alpha$, тогда получим уравнение $\frac{\alpha}{1}=\frac{1-2\alpha}{1-\alpha}$ и $\alpha$ = 0.382. Этот метод на каждом шаге сокращает промежуток на 38.2% независимо от вида унимодальной функции.

3) Метод последовательной параболической интерполяции

На каждом шаге ориентируется на 3 точки: $x_{k-2}, x_{k-1}, x_k$, пот ним строится парабола, вершина определяет значение $x_{k+1}$, которое заменяет $x_{k+2}$

Сочетание методов последовательной параболической интерполяции и золотого сечения демонстрирует процедура-функция FMIN(A, B, F, EPS)

Используется программа FMIN для минимизации F(p). Последовательность вызова:
MAIN → FMIN → F(p) → RKF45 → f(t, x, p).
&

\textbf{Минимизация функции многих переменных, связь с задачей:}
\(f(x) = f(x^{(1)}, \cdots, x^{(m)})\) - функция многих переменных
Если функция непрерывна вместе со своими частными производными, то задача минимизации может быть сведена к решению системы нелинейных уравнений
(1) $grad f(x) = (\frac{\partial f}{\partial x^{(1)}} , \frac{\partial f}{\partial x^{(2)}}, \cdots , \frac{\partial f}{\partial x^{(m)}})^T$ = 0 задающих нулевое значение компонентам вектора градиента f(x). С другой стороны, пусть задана некоторая система нелинейных уравнений
(2) \(\phi_k(x) = \phi(x^{(1)}, \cdots, x^{(m)}) = 0\),   \(k = 1, 2, ..., m\)
Вместо ее решения можно минимизировать некоторую функцию
(3) $f(x) = \sum_{k = 1}^m \phi_k^2(x)$
Ее минимум = 0 и достигается при тех же x, которые удовлетворяют (2) (МОГУТ ПОЯВИТСЯ ПОСТОРОННИЕ ЛОКАЛЬНЫЕ ЭКСТРЕМУМЫ). Минимизация (3) часто проводится специальными методами, учитывающими особый вид минимизируемой функции.

&

\textbf{Понятие линии уровня минимизируемой функции. Примеры методов минимизации функции многих переменных:}
$x^*$ - точка минимума \(f(x)\)
$x_n$ - n-е приближение к \(x*\)
Поверхность уровня (линия уровня) - множество значений вектора x, удовлетворяющих уравнению $f(x) = c_n$
$x_{n+1} = x_n + hv_n$ - новое приближение
Различные методы отличаются друг от друга выбором $v_n$ и h.

1) Метод покоординатного спуска
В качестве $v_n$ выбраны компоненты вектора x.

Обычный (циклический) покоординатный спуск:
Все компоненты $x^{(k)}$  назначены в строго определенном порядке.

Случайный покоординатный спуск:
Выбор очередной компоненты носит случайный характер.

Релаксационный покоординатный спуск:
На каждом шаге в качестве $v_n$ выбирается та компонента, в направлении которой f(x) убывает быстрее всего.

2) Метод градиентного спуска
В качестве $v_n$ выбирается вектор антиградиента.
Использование оптимального значения h: $f(x_{n+1}) = min_{(h>0)}f(x_n - h*grad*f(x_n))$ приводит к методу наискорейшего градиентного спуска.

3) Метод Ньютона
Если имеется хорошее начальное приближение к x*.
матрица Гессе минимизируемой функции f(x): $H=\frac{\partial^2f}{\partial x^2}$
Сам метод выглядит след. образом: $H(x_{n+1}-x_n)=-grad*f(x_n)$
\\\hline


\textbf{Свойства симметрических матриц. Квадратичная функция и ее линии уровня:}
Пусть элементы матрицы A вещественные, а матрица симметрическая $A=A^T$.

1) Свойство 1 - $(Ax, y) = (x, Ay)$
2) Свойство 2 - Собственные значения и собственные векторы симметрической матрицы вещественны.
3) Свойство 3 - Собственные векторы симметрической матрицы, отвечающие различным собственным значениям, ортогональны.
4) Свойство 4 - Диагонализация симметрической матрицы $\Lambda = U^{-1}AU$, $A=U\Lambda U^{-1}$.

Функция многих переменных может быть описана квадратичной функцией:
$f(x)=\frac{1}{2}(Ax,x)+(b,x)$

Линии уровня квадратичной функции представляют собой эллипсоиды, произвольным образом ориентированные в пространстве. Направления осей определяются собственными векторами матрицы Гессе, а их размеры пропорциональны. Если матрица А плохо обусловлена, то линии уровня имеют сильно вытянутый характер и покоординатный спуск становится крайне неэффективным.
&

\textbf{Понятие овражной функции, ее примеры для функции трех переменных:}
Овражная функция — это функция, график которой содержит узкие, вытянутые вдоль одной оси долины (овраги), то есть в одном направлении градиент значительно меньше, чем в другом.
Примером их в трехмерном случае будут следующие два варианта:

a) \( \lambda_{1} \geq \lambda_{2} \gg \lambda_{3} \) – эллипсоид похож на иглу

b) \( \lambda_{1} \gg \lambda_{2} \geq \lambda_{3} \) – эллипсоид похож на диск

\textbf{МЕТОД обобщенного покоординатного спуска}
вычисляем собственные векторы матрицы \(A\). переходим от координат вектора \(x\) к координатам вектора \(z\) и в этих осях применяемый покоординатный спуск.
Если  минимизируемая ф-ия не квадратичная, то определяем матрицу Гессе и ее собственные векторы, минимизируем ф-ию вдоль этих векторов, пока не произойдет значительное замедление сходимости. Затем вычисляем матрицу Гессе в новой точке и тд.

\textbf{дифференциальные уравнения линии спуска}
Пусть все компоненты вектора \(x\) зависят от некоторой переменной \(\tau\) и являются решением следующей системы дифференциальных уравнений: \(\frac{dx(\tau)}{d\tau} = -\text{grad}(x)\)
Производная: \(\frac{dx(\tau)}{d\tau} = \sum_{k=1}^{m}\frac{f}{x^{(k)}}\frac{dx^{(k)}}{d} = (\text{grad}(x), \frac{dx}{d\tau}) = -(\text{grad}f(x), \text{grad}f(x)) < 0\) , с ростом \(\tau\) эта функция непрерывно убывает и достигает минимума в точке, отвечающей нулю ее градиента
&

\textbf{Комплексы параметров в задачах минимизации, методика их выявления.}
В случае если аналитические ф-лы для \(f(x)\) отсутствуют и комплексы явно не определяются, важно:

a) выяснить какие параметры образуют комплексы

b) понять каков минимальный объем доп информации для однозначного определения всех \(x^{(k)}\).

\textbf{Методика выявления комплексов параметров}
Минимизируя \(f(x)\), находим вектор \(x^{*}\), обеспечивающий минимум \(f(x^{*}) = f^{*}\). Если образовываются комплексы, то значению \(f^{*}\) отвечает некоторое множество параметров \(x\) , отражающих эти комплексы. Изменяем один из параметров \(x^{(k)}\) от найденной \(x^{*}\) на
некоторую величину, и фиксируем значение. Проводим минимизацию за счет остальных параметров. Если в точке \(x^{**}\) получается \(f(x^{**}) \approx f^{*}\), то параметр входит в комплексы. Так же в комплексы входят ВСЕ параметры, значения которых различны в \(x^{*}\) и в \(x^{**}\).
После изменяем и фиксируем остальные компоненты вектора \(x\). Когда, определен перечень параметров, входящих в комплексы, \(x^{(k)}\) понижают на единицу (например фиксируют один из параметров) и повторяют процедуру в пространстве меньшей размерности.


\\\hline

\textbf{Барьерные и штрафные функции в задачах условной минимизации:}

Условная минимизация - минимум достигается на границе разрешенной области.
Найти минимум ф-ции \(f(x)\) при ограничениях на вектор $x:$  \(c_k(x) \geq 0, k=1, 2,\cdots , s \tag{17} \)

Вместо \(f(x)\) будем несколько раз минимизировать функцию $F(x):$ \(F(x)=f(x)+T(x, \alpha)  \tag{18} \) , где \(T(x, \alpha)\) – функция, выбираемая специальным образом; \(\alpha\) – параметр, принимающий последовательно уменьшающиеся значения. При этом минимум \(F(x)\) – находится при фиксированном \(\alpha\), результат является начальным приближением для решения при новом \(\alpha\).

\textbf{Метод барьерных функций}
Барьерная ф-ция: \(T(x, \alpha) = \sum_{k=1}^{S}\frac{1}{c_{k}}(x)  \tag{19}\).
Выбираем начальное приближение \(x_{0}\) внутри области возможного изменения \(x\), когда все условия (17) выполняются. С ростом числа итераций \(n, x_{n}\) приближается к границе области, где \(c_{k}(x)=0\), и ф-ция \(T(x, \alpha)\) резко возрастает, образуя «барьер». Полученную точку минимума используем как начальную для минимизации \(T(x, \alpha)\) с меньшим \(\alpha\). «Барьер» сохраняется, но его влияние уменьшается, \(F(x)\) становится ближе к \(f(x)\).

\textbf{Метод штрафных функций}
Начальное приближение \(x_{0}\) выбирается ВНЕ области возможного изменения \(x\), т.е. когда одно или более условий (17) нарушаются. Штрафная ф-ция: \(T(x, \alpha) = \frac{1}{\alpha} \sum_{k=1}^{s}[min(0, c_{k}(x))]^2\). - те же значения, что и ранее. Внутри области допустимого изменения \(x\) штрафная функция равна нулю и \(F(x)\) совпадает с \(f(x)\). При нахождении \(x_{n}\) вне области к минимизируемой функции \(f(x)\) добавляется «штраф» \(T(x, \alpha)\) значение которого возрастает по мере уменьшения \(\alpha . x_{n}\). извне стремится к границе области.

&

\textbf{Понятие некорректно поставленной задачи, примеры таких задач:}

Введем область «исходных данных» U с элементами $u \in U$ и область «решений» Z с элементами $z \in Z$ . При этом решение z находим по исходным данным u в соответствии с некоторым правилом $z = R(u)$. Метрические пространства U и Z могут иметь самую различную природу, определяемую постановкой задачи, поэтому для каждого из них введем расстояния между элементами $\rho_U(u_1, u_2)$ , и $\rho_Z(z_1,z_2)$ ; $u_1, u_2 \in U ; z_1, z_2 \in Z$ ;  . Пусть каждому элементу $u \in U$ отвечает единственное решение $z=R(u)$ ; $z \in Z$ .
Определение 1
Задача определения решения $z=R(u)$ из пространства Z по исходным данным $u \in U$ называется *устойчивой* на пространствах $(Z, U)$, если
$(\forall\epsilon>0)(\exists\delta(\epsilon)>0)(\rho_U(u_1, u_2) \leq \delta(\epsilon) \Rightarrow \rho_Z(z_1, z_2)\leq\epsilon)$
где $z_1=R(u_1), z_2=R(u_2), u_1, u_2 \in U, z_1, z_2 \in Z$

Определение 2
Задача определения решения \(z\) из пространства \(Z\) по исходным данным $u\in U$ называется корректно поставленной на метрических пространствах $(Z,U)$ если выполняются следующие два условия
1) для всякого элемента $u \in U$ существует определяемое однозначно решение $z\in Z;$
2) задача устойчива на пространствах \((Z, U)\).
Нарушение любого из этих требований делает задачу некорректно поставленной.

Пример 1: \(Ax=b\), в случае \(det(A) = 0\) или для прямоугольной матрицы нарушаем первое условие. Что бы сделать его корректным нужна доп информация, например: выбирать решение с наименьшей нормой.

Пример 2: Решение систем линейных дифференциальных уравнений \((x’=Ax)\), когда одно или несколько собственных значений матрицы A имеет положительную вещественную часть \(Re\lambda_k >0\). Нарушено второе условие (устойчивость). Близость решений для \(t \rightarrow \infty\) не удается обеспечить даже при очень малом отличии в начальных условиях.

Пример 3: обратная задача – по заданному решению \(z \in Z\)  нужно восстановить исходные данные \(u \in U\), т.е. \(u=R*(z)\). Весьма популярная задача – подобрать параметры модели так, чтобы результаты моделирования были максимально близки к экспериментальным данным.

&
\textbf{Устойчивость задачи собственных значений и векторов:}

Проблема собственных значений подразумевает отыскание чисел \(\lambda_1, \lambda_2, \cdots, \lambda_m \) (в общем случае комплексных) и соответствующих ненулевых векторов \(u_1, u_2, \ldots, u_m\) удовлетваряющих уравнению \(Au=\lambda u\) где  \(A\) - матрица размера  \(m \times m \). Все методы делятся на 2 группы «полной проблемы собственных значений» (ищем все собственные числа матрицы) и «частичной проблемы собственных значений» (ищем только некоторые \( \lambda_k \)).
\textbf{Устойчивость проблемы собственных значений}
Среди рассматриваемых собственных значений нет кратных. Сохраняя первичные обозначения определим \( \lambda_k \) и \(u_k\) матрицы \(A\)
(\( A u_k= \lambda_k u_k \)), a \(v_k\) как собственные векторы транспонированной матрицы \(A^T\)
\(A^T v_k= \lambda_k v_k \) или \(v_k^T A = \lambda_k v_k^T\).

Умножим равенство \(v_i^T A = \lambda_i v_i^T\) справа на u_k a равенство \(Au_k = \lambda_k u_k\) слева на \(v_i^T\) и вычтем одно из другого:
\( 0 = (\lambda_i - \lambda_k)v_i^T u_k \)  \(v_i^T u_k = 0\) при \(\lambda_i \neq \lambda_k\)
Если векторы имеюь вещественные элементы, то \(v_i^T u_k = u_k^T v_i = (u_k,v_i) =0 \text{ для } i \neq k\)
каждый вектор из одного ряда ортогонален любому вектору из другого, за исключением вектора с тем же индексом. Такое свойство называется биортогональностью, а векторы биортогональными.
удобно считать, что \( v_k^T u_k = 1 \)

Рассмотрим малое изменение матрицы \(\Delta A\). Изменения собственного значения \(\Delta\lambda_k\) и собственного вектора \(\Delta u_k\) с точностью до членов второго порядка малости будут удовлетворять линеаризованному уравнению
\(\Delta Au_k + A\Delta u_k = \Delta\lambda_k u_k + \lambda_k\Delta u_k\)

Путем преобразований получаем:
\(|\Delta\lambda_k|\leq \frac{||\Delta A||\cdot||v_k||\cdot||u_k||}{|v_k^Tu_k|} = C_k||\Delta A|| \quad C_k = \frac{||v_k||\cdot||u_k||}{|v_k^Tu_k|} = \frac{1}{\cos(\varphi_k)}\)
\(C_k\) - коэффициент перекоса, \(\varphi_i\) - угол меж собственными векторами
мы получаем оценку сверху погрешности собственного вектора \(\frac{||\Delta u_k||}{||u_k||} \leq ||\Delta A|| \cdot \sum_{j \neq k} \frac{C_j}{|\lambda_k - \lambda_j|} \)
Изменение собственного вектора аналогичным образом зависит уже от всех коэффициентов перекоса.
наличие близких друг к другу собственных значений может существенно повлиять на изменение собственных векторов
\\\hline

\end{tabular}


\newpage

\begin{center}
  {\large\bf Шпаргалка кр1}
\end{center}

\begin{tabular}{|p{0.315\textwidth}|p{0.315\textwidth}|p{0.315\textwidth}|} % Adjust column widths
    \hline
    \\\hline


    \textbf{Прямой и обратный степенные методы:}
    \textbf{Прямой:} Степенной метод (или прямые итерации) преднозначен для нахождения максимального по модулю собственного значения и соответствующего собственного вектора. Пусть все собственные значения \(\lambda_k\) матрицы А различны и упорядочены след образом \(|\lambda_1| \geq |\lambda_2| \geq \cdots \geq |\lambda_m|\).
    Начальный вектор \(x_0\) разложим по собственным векторам матрицы \(A\): \(x_0 = \alpha_1 u_1 + \alpha_2 u_2 + \cdots + \alpha_m u_m\)
    Предположим \(\alpha_1 \neq 0\) В ходе применения степеного метода строится последовательность векторов \(x_k\)
    \(x_{k+1} = Ax_k \quad x_k = A^k x_0 = \alpha_1 \lambda_1^k u_1 + \cdots + \alpha_m \lambda_m^k u_m  \)
    Для наименьшего угла \(\varphi\) между векторами \(x_k\) и \(u_1\) получаем
    \(\cos \varphi = \frac{|(x_k, u_1)|}{||x_k|| \cdot ||u_1||}\) Так как \(|\lambda_2 / \lambda_1|<1\), \(\cos\varphi \rightarrow 1\) при \(k\rightarrow\infty\) . Иными словами, вектор \(x_k\) сходится к вектору \(u_1\) , отвечающему собственному значению \(\lambda_1\) . При больших значениях k имеет место приближенное равенство
    \(x_{k+1} = A^{k+1}x_0 \approx \lambda_1^{k+1}u_1 \approx \lambda_1 x_k \)
    Это позволяет вычислять приближенно \(\lambda_1\) как отношение каких-либо компонент векторов \(x_{k+1}\) и \(x_{k}\)
    \(   \lambda_1 \approx \frac{(x_{k+1}, b)}{(x_k, b)} \)
    вместо b может быть $x_k$.
    При \(|\lambda_1|<1\) последовательность $x_k$ сходится к нулю, а при \(|\lambda_1|>1\) она неограниченно возрастает. В связи с этим мы нормируем вектор $x_k$ и на (k+1)-м шаге используется вектор \(x_k/||x_k||\).
    Если начальное приближение $x_0$ выбрано неудачно так, что \(\alpha_1=0\) , то теоретически итерационный процесс должен сходиться к $\lambda_2$. Однако на практике по другому
    \textbf{Обратный:} Обратный степенной метод (или обратные итерации) сводится к степенному методу, примененному к обратной матрице $A^{-1}$ . Рабочие формулы метода имеют вид.
    \( Ax_{k+1} = x_k \) \( \lambda_m \approx \frac{(x_k, b)}{(x_{k+1}, b)} \) Максимальное по медулю собственное значение матрицы $A^{-1}$ равно $(1/\lambda_m)$. Аналогично предыдущему
    \( x_{k+1} = A^{-k-1}x_0 \approx(\frac{1}{\lambda_m})^{k+1} *u_1 \approx \frac{1}{\lambda_m} x_k \)
    Если скорость сходимости прямых итераций определяется величиной $|\lambda_2/\lambda_1|$, то в случае обратных итераций эта скорость зависит от $|\lambda_m/\lambda_{m-1}|$. Трудоемкость расчетов резко снижается, если однократно определять LU–разложение матрицы A (например, программой DECOMP) , а затем на каждой итерации решать две системы с треугольными матрицами (например, программой SOLVE). Как и для прямых итераций, рабочие формулы (5) и (6) необходимо дополнить условиями нормировки и использовать вектор \(x_k/||x_k||\).

    &

    \textbf{Свойства преобразования Хаусхолдера:} Рассмотрим произвольную квадратную матрицу \(A\in\mathbb{R}^{m\times m}\) и определим класс «почти треугольных» матриц Хессенберга и покажем, как с помощью преобразований подобия Хаусхолдера привести любую матрицу к такому виду, не изменив её спектр собственных значений. \textbf{1. Матрица Хессенберга} Матрица \(H = (h_{ik})\) называется \textbf{верхней матрицей Хессенберга}, если \(h_{ik} = 0\quad\text{для}\quad i > k+1.\) Внутренне это «почти треугольная» матрица:
    \[
    H =
    \begin{pmatrix}
    * & * & * & \cdots & * \\
    * & * & * & \cdots & * \\
    0 & * & * & \cdots & * \\
    0 & 0 & * & \cdots & * \\
    \vdots & \vdots & \vdots & \ddots & \vdots
    \end{pmatrix},
    \]
    где «$*$» обозначает произвольные (не обязательно нулевые) элементы, а все, что ниже первой поддиагонали, равно нулю. \textbf{2. Преобразование Хаусхолдера} Для вектора \(v\in\mathbb{R}^n\) вводится матрица \(H = I - 2\,\frac{v\,v^T}{v^T v}.\) Часто нормируют \(v\) до единичной длины \(u = \frac{v}{\|v\|} \) , тогда \(H = I - 2\,u\,u^T.\) \textbf{Свойства:} Симметричность: \(H^T = H. \) ; Ортогональность: \(H^T H = I. \) ; Отображение вектора: Пусть \(e_k\) — \(k \) -й базисный вектор, а \(x\in\mathbb{R}^n \) . Если \(v = x + \|x\|\;e_k,\) то \(H\,x = -\|x\|\;e_k.\) \textbf{3. Алгоритм приведения к форме Хессенберга} Пусть \(A\in\mathbb{R}^{m\times m} \) . Для \(i=1,2,\dots,m-2\) выполняем: \begin{enumerate} \item \textbf{Выделение подветктора.} Возьмём \(i \) -й столбец текущей матрицы \(A_{i-1}\) без первых \(i\) элементов: \(x_i = \bigl(a_{i+1,i},\,a_{i+2,i},\,\dots,\,a_{m,i}\bigr)^T \in \mathbb{R}^{m-i}.\) \item \textbf{Построение отражения.} Определим \(v_i = x_i + \|x_i\|\;e_1,\quad u_i = \frac{v_i}{\|v_i\|},\) где \(e_1\in\mathbb{R}^{m-i}\) — первый стандартный вектор. Сформируем матрицу Хаусхолдера \(H_i = I_{\,m-i} - 2\,u_i\,u_i^T.\) \item \textbf{Расширение до размерности \(m \) .} Введём \[
    \[
    U_i = \begin{pmatrix}
        I_i & 0 \\
        0 & H_i
    \end{pmatrix}\in \mathbb{R}^{m\times m}.
    \]


    \item \textbf{Шаг преобразования подобия.} Обновляем матрицу: \(A_i = U_i\,A_{i-1}\,U_i,\quad A_0 = A.\) После этого в \(i \) -м столбце \(A_i\) под главной поддиагональю остаются только нули. \end{enumerate} После \(m-2\) шагов получаем матрицу \(\widetilde A = U_{m-2}\,\cdots\,U_2\,U_1\,A\,U_1\,U_2\,\cdots\,U_{m-2},\) которая имеет верхнюю форму Хессенберга и сохраняет все собственные значения исходной \(A.\) \textbf{Замечание.} Если \(A\) изначально симметрична, то каждый \(U_i\) также сохраняет симметрию, и итоговая \(\widetilde A\) становится трёхдиагональной. &

    &

    \\\hline

\end{tabular}
\end{document}
